% !TeX root = ../main.tex

\chapter{补充材料}
\label{cha:appendix}

\section{论文证据台账说明}
\label{sec:appendix-ledger}

为保证正文中的每个关键数字都可追溯，本文在论文工程中单独维护了证据台账。台账仅服务于内部实验核对，不作为参考文献，不通过 \verb|\cite| 引用。表~\ref{tab:evidence-ledger} 总结了正文核心结论与其对应证据来源。

\begin{table}[htbp]
  \centering
  \small
  \bicaption{核心结论与证据来源对应关系}{Mapping between core claims and internal evidence sources}
  \label{tab:evidence-ledger}
  \begin{tabularx}{\textwidth}{l l X}
    \toprule
    结论层级 & 证据来源 & 正文落点 \\
    \midrule
    一级强结论 & Synth50 三路 clean subset 汇总统计 & 第五章主实验 A \\
    二级机制结论 & 独立 skill 适配实验汇总统计 & 第五章主实验 B \\
    三级观察结论 & Real Category 1 pilot 与 intervention quick pilot & 第五章真实数据局限性分析 \\
    结构性证据 & 结构化 skill manifest 与可读 skill 文档 & 第四章 Skill Refinement 机制 \\
    \bottomrule
  \end{tabularx}
\end{table}

\section{图表来源说明}
\label{sec:appendix-figure-source}

本文所有图表均通过脚本自动生成，并与正文使用同一套证据台账数据。结构图用于表达系统组成与执行流程，数据图用于展示准确率、时延、收益比与失效模式。图注统一采用“本文实验统计”或“本文方法示意图”的表述，不引用本地路径，也不将本地实验文件作为参考文献条目。

\section{引文范围说明}
\label{sec:appendix-citation-scope}

本文参考文献全部来自公开论文、期刊、会议或公开技术报告。为了满足本科论文写作规范并保留完整背景脉络，文末参考文献统一使用公开 BibTeX 条目组织。与此同时，本地 artifact、Markdown 草稿、日志文件以及 CSV/JSON 汇总表均被显式排除在参考文献系统之外。
